% ============================================================
%  PRÉ-PROJETO — PROJETO INTEGRADOR
%  Opus Freelas — Plataforma de Freelancing para Serviços Rurais
%  IFC Campus Concórdia · Técnico em Informática para Internet
%  Stack alinhada ao planejamento GSD (2026): Expo SDK 55 · Hono 4 ·
%  PostgreSQL/PostGIS (Supabase) · Clerk · Mercado Pago (assinatura)
% ============================================================
\documentclass[12pt, a4paper]{article}

% ── Codificação ─────────────────────────────────────────────
\usepackage[utf8]{inputenc}
\usepackage[T1]{fontenc}
\usepackage[brazil]{babel}

% ── Layout — margens ABNT ───────────────────────────────────
\usepackage[top=3cm, bottom=2.5cm, left=3cm, right=2cm,
            includeheadfoot]{geometry}

% ── Tipografia e espaçamento ────────────────────────────────
\usepackage{setspace}
\usepackage{indentfirst}
\setstretch{1.5}
\setlength{\parindent}{1.25cm}
\setlength{\parskip}{0pt}

% ── Cores ───────────────────────────────────────────────────
\usepackage[dvipsnames, table]{xcolor}
\definecolor{ifcGreen}{RGB}{0, 133, 63}
\definecolor{ifcDark}{RGB}{20, 50, 20}
\definecolor{rowGray}{RGB}{245, 245, 245}
\definecolor{codeGray}{RGB}{245, 247, 250}
\definecolor{stackBlue}{RGB}{30, 80, 160}

% ── Cabeçalho e rodapé — SOLUÇÃO ROBUSTA ────────────────────
% Usamos \raisebox para conter o logo dentro da linha de base
% e evitar que o número de página "escape" da área imprimível.
\usepackage{fancyhdr}
\usepackage{graphicx}

\setlength{\headheight}{46pt}   % altura reservada para o header
\setlength{\headsep}{12pt}      % espaço entre header e texto

\pagestyle{fancy}
\fancyhf{}                      % zera tudo

% Logo contido: raisebox garante que não invade a margem
\fancyhead[L]{%
  \raisebox{-0.3\height}{\includegraphics[height=0.9cm]{logo_ifc.png}}%
}
\fancyhead[R]{%
  \begin{tabular}[b]{r}
    \footnotesize\bfseries\color{ifcDark} Projeto Integrador \\[-2pt]
    \footnotesize\color{ifcDark} IFC --- Campus Concórdia
  \end{tabular}%
}
% Rodapé: número centralizado, dentro das margens
\fancyfoot[C]{\small\thepage}
\fancyfoot[L]{}
\fancyfoot[R]{}

\renewcommand{\headrulewidth}{0.6pt}
\renewcommand{\headrule}{%
  \hbox to\headwidth{\color{ifcGreen}\leaders\hrule height 0.6pt\hfill}}
\renewcommand{\footrulewidth}{0pt}

% Páginas "plain" (sumário, 1ª pág de seção)
\fancypagestyle{plain}{%
  \fancyhf{}
  \fancyfoot[C]{\small\thepage}
  \renewcommand{\headrulewidth}{0pt}
  \renewcommand{\footrulewidth}{0pt}
}

% ── Seções ──────────────────────────────────────────────────
\usepackage{titlesec}
\titleformat{\section}
  {\normalfont\large\bfseries\color{ifcDark}}
  {\thesection}{1em}{}[\color{ifcGreen}\titlerule]
\titleformat{\subsection}
  {\normalfont\normalsize\bfseries\color{ifcDark}}
  {\thesubsection}{1em}{}
\titleformat{\subsubsection}
  {\normalfont\normalsize\itshape\color{ifcDark}}
  {\thesubsubsection}{1em}{}
\titlespacing*{\section}   {0pt}{20pt}{10pt}
\titlespacing*{\subsection}{0pt}{14pt}{6pt}
\titlespacing*{\subsubsection}{0pt}{10pt}{4pt}

% ── Hyperlinks ──────────────────────────────────────────────
\usepackage{hyperref}
\hypersetup{
  colorlinks = true,
  linkcolor  = ifcGreen,
  urlcolor   = ifcGreen,
  citecolor  = ifcDark,
  pdftitle   = {Opus Freelas — Pré-Projeto de Projeto Integrador},
  pdfauthor  = {Murilo Jochkeck; Davi Gusso; Mikael Costa; Joao Bordin; Emily Lermen}
}

% ── Tabelas ─────────────────────────────────────────────────
\usepackage{booktabs}
\usepackage{array}
\usepackage{tabularx}
\usepackage{longtable}
\usepackage{multirow}
\usepackage{colortbl}
\usepackage{float}
\usepackage{caption}
\captionsetup{font=small, labelfont=bf}
\renewcommand{\arraystretch}{1.35}

% ── Listas ──────────────────────────────────────────────────
\usepackage{enumitem}
\setlist[itemize]{leftmargin=1.5cm, itemsep=2pt, topsep=4pt}
\setlist[enumerate]{leftmargin=1.5cm, itemsep=2pt, topsep=4pt}

% ── Caixas (tcolorbox) ──────────────────────────────────────
\usepackage{tcolorbox}
\tcbuselibrary{skins, breakable}

% Caixa de tecnologia / IA
\newtcolorbox{techbox}[2][]{
  title     = #2,
  colback   = ifcGreen!4,
  colframe  = ifcGreen!55!ifcDark,
  fonttitle = \bfseries\small\color{ifcDark},
  boxrule   = 0.7pt,
  arc       = 3pt,
  left=8pt, right=8pt, top=5pt, bottom=5pt,
  breakable, #1
}

% Caixa de aviso / destaque
\newtcolorbox{notebox}{
  colback   = stackBlue!6,
  colframe  = stackBlue!50,
  boxrule   = 0.7pt,
  arc       = 3pt,
  left=8pt, right=8pt, top=5pt, bottom=5pt,
  breakable
}

% ── Código inline ───────────────────────────────────────────
\usepackage{amsmath}
\newcommand{\code}[1]{%
  \texttt{\small\colorbox{codeGray}{\vphantom{Ag}#1}}}

% ── Referências ABNT ────────────────────────────────────────
\usepackage[alf]{abntex2cite}

% ============================================================
\begin{document}

% ============================================================
%  CAPA
% ============================================================
\begin{titlepage}
  \thispagestyle{empty}
  \begin{center}
    \includegraphics[width=0.42\textwidth]{logo_ifc.png}\\[0.7cm]
    {\large\bfseries\color{ifcDark}
      Instituto Federal Catarinense --- Campus Concórdia}\\[0.2cm]
    {\normalsize Curso Técnico Integrado em Informática para Internet}\\[0.2cm]
    {\color{ifcGreen}\rule{0.82\textwidth}{1.5pt}}\\[0.8cm]

    {\normalsize Pré-Projeto de Projeto Integrador}\\[0.8cm]

    {\LARGE\bfseries\color{ifcDark} Opus Freelas}\\[0.35cm]
    {\normalsize\itshape\color{ifcDark}
      (projeto integrador --- marca acadêmica: \textit{Opus Freelas})}\\[0.5cm]
    {\large\color{ifcDark}
      Plataforma Multiplataforma de Freelancing\\
      para Serviços Manuais e Rurais na Região da AMAUC}\\[2.0cm]

    \begin{tabular}{rl}
      \multicolumn{2}{c}{\textbf{Equipe de Desenvolvimento}} \\[0.3cm]
      & Murilo Victor Jochkeck \\
      & Davi Ulisses Gusso \\
      & Mikael Ryan Prediger Costa \\
      & João Pedro Stein Bordin \\
      & Emily Lermen \\
    \end{tabular}\\[1.8cm]

    {\color{ifcGreen}\rule{0.82\textwidth}{1.5pt}}\\[0.5cm]
    {\normalsize Concórdia, SC --- 2026}
  \end{center}
\end{titlepage}

% ============================================================
%  SUMÁRIO
% ============================================================
\tableofcontents
\thispagestyle{plain}
\newpage

% ============================================================
%  RESUMO
% ============================================================
\section*{Resumo}
\addcontentsline{toc}{section}{Resumo}

\noindent O presente pré-projeto descreve o planejamento e o desenvolvimento
da plataforma \textbf{Opus Freelas} (referida no contexto acadêmico também
como \textit{Opus Freelas}), voltada à oferta e contratação de serviços
manuais e rurais na região da Associação dos Municípios do Alto Uruguai
Catarinense (AMAUC). A solução conecta prestadores --- capineiros,
roçadores, operadores de máquinas agrícolas, diaristas e autônomos em
geral --- a contratantes locais, reduzindo a informalidade da descoberta
de serviços e aumentando confiança por meio de perfil, localidade,
verificação básica e reputação.

O escopo da primeira versão (\textit{minimum viable product}), definido de
forma explícita no planejamento do produto, \textbf{prioriza o contratante}
(publicação e descoberta de demanda) e adota \textbf{fechamento do serviço
fora da plataforma} (telefone/WhatsApp), sem intermediação financeira entre
as partes nessa etapa. A monetização inicial prevista é \textbf{assinatura
do prestador} (modelo SaaS), com integração ao ecossistema brasileiro via
\textbf{Mercado Pago}. A arquitetura combina aplicação \textbf{multiplataforma}
com \textbf{Expo} (React Native, SDK alinhado à documentação 2026), API em
\textbf{Node.js} com \textbf{Hono}, dados em \textbf{PostgreSQL} com
\textbf{PostGIS} (infraestrutura gerenciada via \textbf{Supabase}),
autenticação com \textbf{Clerk} (fluxo por telefone com OTP, papéis de
contratante e prestador na mesma conta) e observabilidade com padrões
modernos (logs, métricas, \textit{tracing}).

O desenvolvimento segue metodologia de entregas iterativas com apoio de
ferramentas de IA e gestão de trabalho estruturada pelo fluxo
\textbf{GSD} (\textit{Get Shit Done}), com repositório versionado,
artefatos de requisitos, roadmap em fases e rastreabilidade entre escopo
e implementação.

\medskip
\noindent\textbf{Palavras-chave:} freelancing; serviços rurais; AMAUC;
Expo; React Native; Hono; PostgreSQL; PostGIS; Supabase; Clerk;
Mercado Pago; LGPD.

\newpage

% ============================================================
%  1. INTRODUÇÃO
% ============================================================
\section{Introdução}

O avanço das tecnologias digitais impulsionou o surgimento da chamada
\textit{economia de plataformas}, modelo no qual intermediários digitais
conectam ofertantes e demandantes de serviços de forma eficiente e
transparente, eliminando barreiras geográficas e reduzindo custos de
transação \cite{workana, getninjas}. Plataformas como Workana,
99Freelas, GetNinjas e Uber consolidaram esse modelo em centros urbanos e
no mercado de serviços digitais.

Contudo, o meio rural permanece amplamente excluído desse fenômeno. Na
região da AMAUC --- que concentra forte vocação para a agricultura
familiar, criação de suínos e aves e o cultivo de soja e milho ---,
a contratação de serviços como capina, roçada, operação de máquinas,
plantio, colheita, cercamento e limpeza de terrenos ainda ocorre de
forma inteiramente informal: por indicação boca a boca, sem contrato,
sem garantia de pagamento e sem histórico de reputação para nenhuma das
partes \cite{ibge2019, sebrae2022}.

Essa lacuna representa simultaneamente um problema socioeconômico e uma
oportunidade concreta de aplicação tecnológica. A plataforma \textit{Opus Freelas}
propõe uma solução multiplataforma --- com prioridade para dispositivos móveis
e, quando aplicável, alvo web --- que formaliza
digitalmente essas relações de trabalho, oferecendo aos prestadores
visibilidade e renda consistente, e aos contratantes rapidez, segurança
e rastreabilidade na contratação. O projeto é desenvolvido como projeto
integrador do Curso Técnico Integrado em Informática para Internet do
IFC Campus Concórdia.

% ============================================================
%  2. OBJETIVOS
% ============================================================
\section{Objetivos}

\subsection{Objetivo Geral}

Projetar, desenvolver e entregar a plataforma \textbf{Opus Freelas}:
solução multiplataforma para contratação de serviços manuais e rurais na
AMAUC, com \textbf{foco inicial no contratante}, autenticação e perfis
consistentes, descoberta por localidade e tipo de serviço, mecanismos de
confiança (verificação, portfólio, avaliações) e \textbf{fechamento do
negócio fora da plataforma} na primeira versão, conforme decisões de produto
documentadas no planejamento (GSD).

\subsection{Objetivos Específicos}

\begin{itemize}
  \item Formalizar requisitos funcionais e não funcionais com identificadores
        rastreáveis, priorização de MVP e registro explícito do que permanece
        fora do escopo da primeira versão (ex.: contratos digitais formais,
        intermediação financeira entre as partes);
  \item Modelar e implementar o domínio em \textbf{PostgreSQL} com
        \textbf{PostGIS} quando necessário para raio e município, com
        políticas de \textbf{Row Level Security (RLS)} alinhadas ao modelo
        de identidade (integração Clerk + Supabase);
  \item Implementar o cliente multiplataforma com \textbf{Expo}
        (React Native), navegação previsível (ex.: \textit{expo-router}),
        cache de dados (\textit{TanStack Query}) e validação compartilhada
        (\textbf{Zod});
  \item Implementar API em \textbf{Node.js} com \textbf{Hono}, rotas em
        estilo RPC para identidade e regras de negócio, validação na borda
        e testes automatizados na API (\textbf{Vitest});
  \item Configurar autenticação com \textbf{Clerk} (telefone com OTP,
        sessão conforme política definida, recuperação com apoio de e-mail
        de backup), com autorização por papéis de contratante e prestador
        na mesma conta quando aplicável;
  \item Integrar \textbf{Mercado Pago} para \textbf{assinatura do
        prestador} (visibilidade e benefícios), sem depender de
        intermediação de pagamento entre contratante e prestador no V1;
  \item Estabelecer observabilidade (registros estruturados, erros,
        métricas e \textit{tracing}) e integração futura com serviços de
        suporte (ex.: \textbf{Firebase} para notificações), sem duplicar
        diretório de usuários;
  \item Estruturar o fluxo de desenvolvimento com Git/GitHub, integração
        contínua em ambiente local e gestão de fases conforme metodologia
        \textbf{GSD};
  \item Utilizar ferramentas de IA como copiloto de desenvolvimento e
        pesquisa, com revisão humana obrigatória;
  \item Validar a solução com testes automatizados onde aplicável e
        verificações manuais ou instrumentadas para fluxos que dependem de
        SMS, dispositivos reais e políticas de loja.
\end{itemize}

% ============================================================
%  3. JUSTIFICATIVA
% ============================================================
\section{Justificativa}

O projeto \textbf{Opus Freelas} justifica-se pela confluência de três fatores:
lacuna de mercado real, alinhamento educacional e escolha tecnológica
estratégica.

\textbf{Lacuna de mercado.} A AMAUC é composta por municípios
predominantemente rurais onde a contratação de serviços manuais ocorre
sem suporte tecnológico ou garantias legais. Dados do IBGE apontam
elevado índice de trabalhadores autônomos informais no setor agrícola do
oeste catarinense \cite{ibge2019}, e o SEBRAE indica que a digitalização
de serviços é o principal vetor de inclusão produtiva nesse perfil de
região \cite{sebrae2022}. Uma plataforma que centralize oferta, demanda
e sinais de confiança (localidade, verificação, reputação) contribui
para reduzir assimetria de informação e informalidade na \textit{descoberta}
do serviço, aumentando a renda dos prestadores e reduzindo o risco percebido
pelos contratantes --- mesmo quando o pagamento entre as partes permanece
fora da plataforma na primeira versão, como decisão explícita de produto.

\textbf{Alinhamento educacional.} O projeto integra, em uma única
aplicação funcional, as competências centrais do Curso Técnico Integrado
em Informática para Internet: desenvolvimento mobile multiplataforma
(React Native via Expo), API e domínio em TypeScript (Node.js, Hono),
banco de dados relacional e geoespacial (PostgreSQL, PostGIS),
design de interfaces (Figma), versionamento (Git/GitHub) e segurança
da informação (RLS, JWT via Clerk, LGPD) \cite{sommerville2019, pressman2016}.

\textbf{Escolha tecnológica.} A stack definida no planejamento --- cliente
\textit{Expo} (ecossistema React Native), API \textit{Hono} em Node.js,
PostgreSQL com PostGIS (via Supabase gerenciado), autenticação Clerk,
integração Mercado Pago para assinatura e observabilidade moderna ---
alinha produtividade de equipe enxuta, tipagem em TypeScript e caminhos
oficiais documentados para 2026, reduzindo retrabalho e ambiguidade entre
protótipo e produção \cite{react2024, postgresql2024}.

% ============================================================
%  4. ESCOPO DO SISTEMA
% ============================================================
\section{Escopo do Sistema}

\subsection{Perfis de Usuário}

A plataforma opera com dois papéis --- \textbf{contratante} e
\textbf{prestador} --- que podem coexistir na \textbf{mesma conta}, com
autorização explícita por papel nas operações sensíveis (evitando vazamento
de permissões apenas na interface).

\begin{itemize}
  \item \textbf{Contratante:} perfil priorizado no MVP; publica demandas,
        busca prestadores por tipo de serviço e localidade, registra
        status da demanda e avalia após conclusão; o fechamento comercial
        ocorre preferencialmente fora da plataforma (contato direto).
  \item \textbf{Prestador de serviço:} mantém perfil com localidade,
        raio de atendimento, portfólio e verificação básica; pode aderir a
        plano de assinatura para maior visibilidade; responde a demandas e
        acumula reputação.
\end{itemize}

\subsection{Módulos do Sistema}

O Quadro~\ref{tab:modulos} organiza os módulos conforme o escopo da
\textbf{primeira versão} e entregas futuras, alinhado ao planejamento por
fases (autenticação e base; demandas; descoberta; confiança; conexão e
reputação; assinatura; endurecimento).

\begin{table}[H]
  \centering
  \caption{Módulos do sistema e prioridade de entrega (alinhado ao roadmap)}
  \label{tab:modulos}
  \begin{tabularx}{\textwidth}{%
    >{\raggedright\arraybackslash}p{4.2cm}
    >{\centering\arraybackslash}p{1.8cm}
    >{\raggedright\arraybackslash}X}
  \toprule
  \rowcolor{ifcGreen!20}
  \textbf{Módulo} & \textbf{Prioridade} & \textbf{Descrição resumida} \\
  \midrule
  Identidade e acesso       & MVP & Clerk (OTP telefone), sessão, papéis contratante/prestador, API Hono + RLS \\
  \rowcolor{rowGray}
  Demandas                  & MVP & Publicação, edição e encerramento; visibilidade por município/raio \\
  Descoberta                & MVP & Busca e filtros por serviço e localidade; estados vazios claros \\
  \rowcolor{rowGray}
  Confiança (prestador)     & MVP & Verificação básica, portfólio com mídia \\
  Conexão e reputação       & MVP & Handoff de contato (telefone/WhatsApp), status da demanda, avaliações \\
  \rowcolor{rowGray}
  Assinatura (prestador)    & MVP & Mercado Pago --- recorrência/assinatura; regras de visibilidade \\
  Geolocalização            & MVP & Consultas por município e raio (PostGIS); mapas conforme integração escolhida \\
  \rowcolor{rowGray}
  Notificações Push         & MVP/Futuro & Firebase (FCM) como canal de suporte; evolução conforme necessidade \\
  Pagamento entre partes    & Fora do V1 & Intermediação financeira e \textit{escrow} --- explicitamente excluídos na v1 \\
  \rowcolor{rowGray}
  Chat in-app completo      & Futuro & Mensageria rica pode ser adicionada após liquidez e fechamento off-platform \\
  Painel administrativo     & Futuro & Moderação avançada, disputas formais, relatórios \\
  \bottomrule
  \end{tabularx}
\end{table}

\subsection{Categorias de Serviço (MVP)}

\begin{itemize}
  \item Serviços Rurais: capina, roçada, carpição, limpeza de pastos,
        coroamento de plantas;
  \item Serviços com Máquinas: trator, roçadeira mecânica,
        retroescavadeira, caminhão basculante;
  \item Construção e Reforma: pedreiro, pintor, servente, assentamento
        de tijolos;
  \item Serviços Agrícolas: plantio, colheita, adubação, irrigação;
  \item Limpeza e Conservação: limpeza de terrenos, poda de árvores,
        remoção de entulho;
  \item Serviços Gerais: carregamento, mudanças, jardinagem.
\end{itemize}

% ============================================================
%  5. STACK TECNOLÓGICA
% ============================================================
\section{Stack Tecnológica}

\subsection{Visão Geral da Arquitetura}

A arquitetura adota \textbf{monorepo} (\textit{pnpm} + Turborepo), com
\textbf{cliente multiplataforma} em \textbf{Expo} (React Native), uma
\textbf{API própria} em \textbf{Node.js} com \textbf{Hono} para regras de
negócio e validação na borda, e \textbf{PostgreSQL} com \textbf{PostGIS}
hospedado via \textbf{Supabase} (migrações, RLS, Storage). A identidade
fica centralizada no \textbf{Clerk} (JWT, OTP por telefone); o app e a API
validam tokens e aplicam autorização por papéis. Essa separação evita
acoplar toda a lógica ao modelo ``somente BaaS'' e mantém o produto
alinhado ao roadmap em sete fases (fundação e identidade; publicação de
demandas; descoberta local; base de confiança; conexão e reputação;
monetização por assinatura; endurecimento). A Tabela~\ref{tab:stack}
resume a stack.

\begin{table}[H]
  \centering
  \caption{Stack tecnológica --- Opus Freelas (planejamento GSD, 2026)}
  \label{tab:stack}
  \begin{tabularx}{\textwidth}{%
    >{\raggedright\arraybackslash}p{2.8cm}
    >{\raggedright\arraybackslash}p{3.8cm}
    >{\raggedright\arraybackslash}X}
  \toprule
  \rowcolor{ifcGreen!20}
  \textbf{Camada} & \textbf{Tecnologia} & \textbf{Bibliotecas / Ferramentas} \\
  \midrule
  Cliente (mobile / web alvo)
    & Expo SDK 55 / React Native 0.83
    & expo-router, TanStack Query, Zod, i18next \\
  \rowcolor{rowGray}
  API e domínio
    & Node.js 22 LTS + Hono 4
    & Drizzle ORM, Vitest, Zod \\
  Dados
    & PostgreSQL 17 + PostGIS (Supabase)
    & Migrações SQL, RLS, índices espaciais \\
  \rowcolor{rowGray}
  Identidade
    & Clerk
    & OTP telefone, JWT, @clerk/clerk-expo \\
  Pagamentos (V1)
    & Mercado Pago
    & Assinatura/recorrência do \textbf{prestador} (SaaS) \\
  \rowcolor{rowGray}
  Mídia
    & Supabase Storage
    & URLs assinadas; buckets com políticas \\
  Geolocalização
    & PostGIS + mapas (conforme integração)
    & Consultas por município/raio; mapa nativo ou web \\
  \rowcolor{rowGray}
  Notificações
    & Firebase (FCM)
    & Canal de suporte; configuração via Expo \\
  Observabilidade / CI
    & GitHub Actions, logs, OpenTelemetry
    & Pino, rastreio de requisições \\
  \rowcolor{rowGray}
  Build loja
    & EAS Build / Submit
    & Perfis dev, preview e produção \\
  Versionamento
    & Git + GitHub
    & Conventional Commits, PRs, proteção de branch \\
  \rowcolor{rowGray}
  Design
    & Figma
    & Design system, wireframes, protótipo navegável \\
  Gestão de trabalho
    & GSD + GitHub Projects
    & Roadmap em fases, requisitos rastreáveis \\
  \bottomrule
  \end{tabularx}
\end{table}

\subsection{Cliente multiplataforma --- Expo (React Native)}

O \textbf{Expo} \cite{react2024} concentra o cliente em React Native com
SDK alinhado à matriz oficial de 2026, permitindo um único código para
Android e iOS (e alvo web quando aplicável). A navegação usa
\textit{expo-router} (rotas previsíveis e \textit{deep links} --- úteis
para compartilhar demandas). O estado servidor/cliente prioriza
\textbf{TanStack Query} para cache e revalidação em redes intermitentes;
\textbf{Zod} compartilha contratos com a API. Telas críticas tratam
estados vazios, erro e carregamento de forma explícita, coerente com o
público-alvo regional.

\subsection{API e domínio --- Hono, Drizzle e testes}

A camada HTTP é implementada com \textbf{Hono} sobre Node.js~22 LTS:
rotas REST/RPC para identidade, demandas, descoberta e assinatura, com
validação na entrada e saída (\textbf{Zod}). \textbf{Drizzle ORM}
acessa o PostgreSQL com SQL explícito --- importante para consultas
PostGIS --- mantendo tipagem forte. Testes automatizados na API usam
\textbf{Vitest}; integração contínua valida regressões antes do
\textit{merge}.

\subsection{Dados e plataforma --- Supabase}

O \textbf{Supabase} \cite{supabase2024} hospeda PostgreSQL gerenciado,
Storage e ferramentas de migração; \textbf{não} substitui a API de
domínio como único ponto de regras de negócio. Destacam-se: extensão
\textbf{PostGIS} para raio e município; \textbf{Row Level Security (RLS)}
como última linha de defesa alinhada a políticas por usuário e papel;
\textbf{Storage} para mídia de portfólio com URLs assinadas. Eventuais
webhooks de cobrança (ex.: Mercado Pago) são tratados na API com
segredo de serviço, sem expor chaves ao cliente.

\subsection{Identidade --- Clerk}

O \textbf{Clerk} fornece autenticação com \textbf{OTP por telefone},
sessão configurável e emissão de \textbf{JWT} consumidos pela API Hono e
pelo app (\texttt{@clerk/clerk-expo}, armazenamento seguro no
dispositivo). Papéis de \textbf{contratante} e \textbf{prestador}
coexistem na mesma conta quando aplicável; a API reforça autorização por
papel para evitar vazamento só pela interface.

\subsection{Banco de Dados --- PostgreSQL com RLS}

A Tabela~\ref{tab:tabelas} descreve entidades alinhadas ao escopo V1
(demanda do contratante, descoberta, conexão sem pagamento entre
partes, assinatura do prestador). Nomes finais seguem o repositório;
a lista fixa conceitos de domínio para o pré-projeto.

\begin{table}[H]
  \centering
  \caption{Modelo de dados --- entidades principais (visão conceitual V1)}
  \label{tab:tabelas}
  \begin{tabularx}{\textwidth}{%
    >{\ttfamily\raggedright\arraybackslash}p{2.8cm}
    >{\raggedright\arraybackslash}X
    >{\raggedright\arraybackslash}p{4.2cm}}
  \toprule
  \rowcolor{ifcGreen!20}
  \textbf{Tabela} & \textbf{Descrição} & \textbf{Colunas-chave} \\
  \midrule
  profiles
    & Perfil de usuário; vínculo a identidade Clerk; papéis e dados públicos
    & id, clerk\_user\_id, papéis, nome, município, foto\_url \\
  \rowcolor{rowGray}
  demands
    & Demandas publicadas pelo contratante (ciclo de vida e visibilidade)
    & id, autor\_id, tipo\_serviço, descrição, status, local, geom \\
  provider\_profile
    & Dados do prestador: raio, categorias, verificação, provas de serviço
    & user\_id, raio\_km, categorias[], verificado\_em \\
  \rowcolor{rowGray}
  reviews
    & Avaliações após conclusão registrada na plataforma
    & id, demand\_id, avaliador\_id, avaliado\_id, nota, comentario \\
  subscription\_events
    & Eventos de assinatura Mercado Pago (prestador); estado para regras de visibilidade
    & id, user\_id, mp\_id, status, periodo \\
  \bottomrule
  \end{tabularx}
\end{table}

\textbf{Estados da demanda (exemplo):} \code{aberta}
$\rightarrow$ \code{em\_contato} $\rightarrow$ \code{concluida}
$|$ \code{cancelada}. O fechamento financeiro entre as partes ocorre
fora da plataforma na V1; a plataforma registra status e reputação.

\subsection{Pagamentos --- Mercado Pago (assinatura do prestador)}

O \textbf{Mercado Pago} \cite{mercadopago2024} é adotado pelo ecossistema
brasileiro (Pix, cartões, hábitos locais). Na \textbf{primeira versão},
o foco de integração é \textbf{assinatura ou recorrência do prestador}
(modelo SaaS da plataforma), com benefícios claros de visibilidade ---
\textbf{sem} custódia (\textit{escrow}) do valor do serviço entre
contratante e prestador. Fluxo típico:

\begin{enumerate}
  \item Prestador adere ao plano via fluxo configurado (checkout ou
        redirecionamento seguro);
  \item Webhooks e estado de assinatura são processados na API (Node.js),
        atualizando permissões e regras de destaque;
  \item Intermediação financeira \textit{business-to-business} entre
        partes do serviço permanece fora do escopo V1 e pode ser
        avaliada em versões futuras.
\end{enumerate}

\subsection{Geolocalização --- PostGIS e mapas}

A extensão \textbf{PostGIS} do PostgreSQL permite armazenar geometrias
(pontos de referência, demandas ou prestadores) e executar consultas por
raio e município diretamente no banco:

\begin{notebox}
\small\ttfamily
SELECT * FROM demands \\
WHERE ST\_DWithin(geom, \\
\quad ST\_MakePoint(-52.02, -27.23)::geography, 30000);
-- exemplo: demandas ou ofertas em até 30 km de Concórdia
\end{notebox}

O mapa no cliente (ex.: \textit{react-native-maps} ou Google Maps,
conforme integração escolhida no Expo) exibe marcadores de forma
responsiva; \textit{clustering} pode ser aplicado quando houver muitos
resultados na mesma área.

\subsection{Ferramentas de IA --- Copiloto de Desenvolvimento}

O desenvolvimento seguirá a metodologia \textbf{vibe coding} assistida
por IA: a equipe descreve a funcionalidade desejada em linguagem natural
e usa as ferramentas abaixo como copiloto, revisando, adaptando e
integrando o código gerado.

\begin{techbox}{Claude (Anthropic) \cite{claude2024} --- Copiloto Principal de Código}
Geração de telas e hooks Expo/React Native com TypeScript; rotas Hono e
middleware Clerk; queries SQL com PostGIS e migrações Drizzle; schemas
Zod compartilhados; revisão e refatoração com explicação dos trade-offs.
\end{techbox}

\vspace{5pt}
\begin{techbox}{ChatGPT (OpenAI) \cite{chatgpt2024} --- Redação e UX Writing}
Textos institucionais da plataforma (FAQs, onboarding, descrições de
categorias em linguagem acessível ao público rural); documentação
técnica (README, comentários de código); brainstorming de funcionalidades
e fluxos de usuário com base em referências de UX.
\end{techbox}

\vspace{5pt}
\begin{techbox}{Gemini (Google) \cite{gemini2024} --- Análise Visual e Pesquisa}
Análise de adequação de imagens de perfil e portfólio enviadas pelos
usuários; pesquisa sobre boas práticas de UX para públicos com baixo
letramento digital; suporte multilíngue para populações com descendência
alemã e italiana presentes na região da AMAUC.
\end{techbox}

\vspace{5pt}
\begin{techbox}{Grok (xAI) \cite{grok2024} --- Contextualização e Tendências}
Pesquisa de tendências atuais no mercado de plataformas de serviços e
gig economy; comparação de abordagens técnicas em tempo real; segunda
opinião sobre decisões de arquitetura e escolhas de biblioteca.
\end{techbox}

\vspace{5pt}
\begin{techbox}{NotebookLM (Google) \cite{notebooklm2024} --- Síntese Bibliográfica}
Organização e síntese do referencial teórico: carregamento de artigos,
leis (LGPD, CLT) e documentos técnicos para geração de resumos e
perguntas diretas sobre o conteúdo. Acelera a revisão bibliográfica e a
redação das seções acadêmicas do projeto.
\end{techbox}

\vspace{5pt}

% ============================================================
%  6. METODOLOGIA DE DESENVOLVIMENTO
% ============================================================
\section{Metodologia de Desenvolvimento}

\subsection{Vibe Coding com GSD}

O processo de desenvolvimento adotará a metodologia
\textbf{vibe coding} sobre a estrutura de gestão
\textbf{GSD (\textit{Get Shit Done})}: foco em entregas concretas e
funcionais a cada ciclo curto, sem cerimônia desnecessária. A cada
semana a equipe define \textbf{3 a 5 tarefas fecháveis} no GitHub
Projects, executa e marca como done antes de planejar a próxima.

Princípios operacionais:

\begin{itemize}
  \item \textbf{Uma tarefa por vez por pessoa:} ninguém abre nova
        tarefa sem fechar a anterior;
  \item \textbf{Done significa deployado:} código que não está em
        \code{develop} não conta como feito;
  \item \textbf{IA é rascunho, humano é revisão:} todo código gerado
        por IA é revisado, testado e entendido pela equipe antes de ir
        para o repositório;
  \item \textbf{Reunião de sync semanal de 30 min:} o que foi feito,
        o que será feito, o que está bloqueado.
\end{itemize}

\subsection{Estratégia de Branches (Git Flow Simplificado)}

\begin{itemize}
  \item \code{main} --- código em produção; protegida, merge somente
        via PR aprovado;
  \item \code{develop} --- integração; toda feature passa aqui antes
        de ir para main;
  \item \code{feature/nome-curto} --- uma feature por branch
        (ex.: \code{feature/busca-por-raio});
  \item \code{fix/descricao-bug} --- correções pontuais.
\end{itemize}

Commits seguem \textbf{Conventional Commits}:
\code{feat:}, \code{fix:}, \code{docs:}, \code{chore:}, \code{style:}.

\subsection{Fluxo Geral de Uso da Plataforma}

O Quadro~\ref{tab:fluxo} descreve o caminho completo de um usuário
desde o primeiro acesso até a conclusão e avaliação de um serviço.

\begin{table}[H]
  \centering
  \caption{Fluxo geral de uso --- Opus Freelas (V1, fechamento externo)}
  \label{tab:fluxo}
  \begin{tabularx}{\textwidth}{%
    >{\centering\bfseries\arraybackslash}p{1.1cm}
    >{\centering\arraybackslash}p{3.2cm}
    >{\raggedright\arraybackslash}X}
  \toprule
  \rowcolor{ifcGreen!20}
  \textbf{Etapa} & \textbf{Ator} & \textbf{Ação} \\
  \midrule
  1  & Visitante          & Explora categorias e contexto; pode criar conta para publicar ou responder \\
  \rowcolor{rowGray}
  2  & Novo usuário       & Autentica-se (Clerk, OTP telefone); habilita papéis contratante e/ou prestador \\
  3  & Contratante        & Publica demanda (tipo, local, urgência) com visibilidade por município/raio \\
  \rowcolor{rowGray}
  4  & Prestador          & Completa perfil, portfólio e verificação básica; opcionalmente assina plano \\
  5  & Contratante        & Busca e filtra prestadores por serviço, município e reputação \\
  \rowcolor{rowGray}
  6  & Contratante        & Abre perfil; solicita revelação de contato ou inicia fluxo de handoff \\
  7  & Ambos              & Negociam e fecham \textbf{fora da plataforma} (telefone/WhatsApp) \\
  \rowcolor{rowGray}
  8  & Ambos              & Atualizam status da demanda na plataforma (ex.: em contato, concluída) \\
  9  & Prestador          & Executa o serviço (fora do escopo técnico da plataforma) \\
  \rowcolor{rowGray}
  10 & Ambos              & Registram conclusão na demanda; habilitam avaliação \\
  11 & Ambos              & Avaliam (1--5 estrelas + comentário), conforme regras anti-abuso \\
  \bottomrule
  \end{tabularx}
\end{table}

\subsection{Design e Prototipagem}

Antes de qualquer código de tela ser escrito, todas as telas serão
prototipadas no \textbf{Figma} \cite{figma2024}. O processo de design
seguirá as etapas:

\begin{enumerate}
  \item \textbf{Design System:} definição de paleta de cores,
        tipografia (Inter), espaçamentos, grid e componentes base
        (botões, cards, inputs, badges);
  \item \textbf{Wireframes de baixa fidelidade:} esboço de todos os
        fluxos principais em escala mobile-first;
  \item \textbf{Protótipo de alta fidelidade:} telas completas com
        estados reais (vazio, carregando, com dados, com erro);
  \item \textbf{Validação:} sessão de teste com ao menos 5 usuários
        do público-alvo da AMAUC antes do início do desenvolvimento.
\end{enumerate}

Os princípios de usabilidade de Nielsen \cite{nielsen1990} norteiam
as decisões de interface, com ênfase em feedback imediato, prevenção
de erros e linguagem acessível para usuários com baixo letramento
digital.

\subsection{Segurança e Conformidade com a LGPD}

\begin{itemize}
  \item \textbf{Row Level Security (RLS):} políticas no PostgreSQL
        garantem que cada usuário só acesse linhas permitidas, mesmo
        quando o cliente usa chaves públicas do Supabase;
  \item \textbf{JWT Clerk (API):} a API Hono valida tokens emitidos pelo
        Clerk; segredos de serviço e webhooks não são expostos ao app;
  \item \textbf{Armazenamento seguro no dispositivo:} sessão e tokens
        sensíveis seguem práticas do \texttt{@clerk/clerk-expo} (ex.\
        SecureStore);
  \item \textbf{Variáveis de ambiente:} chaves de Supabase, Clerk e
        Mercado Pago apenas em CI e hospedagem (nunca no repositório);
  \item \textbf{Conformidade LGPD \cite{lgpd2018}:} coleta mínima de
        dados, política de privacidade acessível, direito de exclusão
        de conta com remoção de dados pessoais;
  \item \textbf{Validação dupla:} dados validados no cliente (Zod) e
        novamente na API (Hono + Zod).
\end{itemize}

\subsection{Testes e Validação}

\begin{itemize}
  \item \textbf{Testes funcionais:} verificação manual dos fluxos
        críticos (cadastro, demanda, descoberta, handoff, status,
        avaliação, assinatura do prestador) no app Expo;
  \item \textbf{Testes de integração:} comunicação entre cliente
        $\leftrightarrow$ API Hono $\leftrightarrow$ PostgreSQL
        (RLS) e webhooks Mercado Pago em ambiente de staging;
  \item \textbf{Testes automatizados:} Vitest na API; fluxos E2E
        prioritários quando a fase de endurecimento estiver ativa;
  \item \textbf{Testes de usabilidade:} sessões com usuários reais
        da AMAUC (agricultores, diaristas, operadores de máquinas)
        para identificar barreiras de adoção;
  \item \textbf{Testes de segurança básicos:} acesso a rotas
        protegidas sem JWT válido, dados malformados na API,
        verificação das políticas RLS no Supabase.
\end{itemize}

% ============================================================
%  7. RESULTADOS ESPERADOS
% ============================================================
\section{Resultados Esperados}

\begin{itemize}
  \item Aplicativo multiplataforma (Expo/React Native) e API (Hono)
        com módulos do MVP integrados a Supabase (dados), Clerk
        (identidade) e Mercado Pago (assinatura do prestador);
  \item Banco PostgreSQL com PostGIS, RLS alinhado ao modelo de
        identidade, diagrama ER e dados de teste;
  \item Fluxo de assinatura do prestador operacional via Mercado Pago,
        com regras de visibilidade refletidas no produto;
  \item Interface adequada a dispositivos móveis e contexto rural,
        validada com usuários reais da AMAUC;
  \item Repositório GitHub em monorepo, branches protegidas e CI
        reproduzível;
  \item Deploy da API em hospedagem compatível com Node.js (ex.\
        Fly.io, Railway ou Render) e binários mobile via EAS Build;
  \item Contribuição concreta para reduzir informalidade na descoberta
        de serviços na AMAUC;
  \item Capacitação prática da equipe em Expo, Hono, Drizzle, Clerk,
        Supabase, Mercado Pago, PostGIS e uso de IA como copiloto.
\end{itemize}

% ============================================================
%  8. CRONOGRAMA DE EXECUÇÃO
% ============================================================
\section{Cronograma de Execução}

\begin{longtable}{%
  >{\centering\arraybackslash}p{0.7cm}
  p{7.2cm}
  p{3.4cm}
  >{\centering\arraybackslash}p{1.4cm}}

  \caption{Cronograma detalhado de execução --- 2026}
  \label{tab:cronograma} \\

  \toprule
  \rowcolor{ifcGreen!20}
  \textbf{\#} & \textbf{Tarefa} & \textbf{Período} & \textbf{Resp.} \\
  \midrule
  \endfirsthead

  \toprule
  \rowcolor{ifcGreen!20}
  \textbf{\#} & \textbf{Tarefa} & \textbf{Período} & \textbf{Resp.} \\
  \midrule
  \endhead

  \bottomrule
  \endfoot

  1  & Elaboração e aprovação do pré-projeto
     & Mar. 2026 & Todos \\
  \rowcolor{rowGray}
  2  & Levantamento de requisitos com usuários reais da AMAUC
     & Mar./Abr. & Todos \\
  3  & Definição do escopo MVP e backlog inicial no GitHub Projects
     & Abr. 2026 & Murilo, Davi \\
  \rowcolor{rowGray}
  4  & Modelagem do banco (diagrama ER) e scripts SQL iniciais
     & Abr. 2026 & Mikael \\
  5  & Configuração do projeto Supabase (tabelas, RLS, Auth, Storage)
     & Abr. 2026 & Mikael, João P. \\
  \rowcolor{rowGray}
  6  & Design System no Figma (cores, tipografia, componentes base)
     & Abr. 2026 & Emily \\
  7  & Wireframes e protótipo de alta fidelidade (mobile-first)
     & Abr./Mai. & Emily, João P. \\
  \rowcolor{rowGray}
  8  & Validação do protótipo com usuários reais da AMAUC
     & Mai. 2026 & Todos \\
  9  & Monorepo (pnpm, Turborepo), pacote compartilhado e contratos Zod
     & Mai. 2026 & Murilo, Mikael \\
  \rowcolor{rowGray}
  10 & API Hono + Drizzle: baseline, rotas iniciais e smoke (Vitest)
     & Mai. 2026 & Mikael \\
  11 & App Expo (SDK 55): expo-router, shell e fluxos de navegação
     & Mai. 2026 & Murilo, Davi \\
  \rowcolor{rowGray}
  12 & Supabase: migrações SQL, RLS, Storage e evidência de deploy local
     & Mai. 2026 & Mikael, João P. \\
  13 & Clerk na API: validação JWT, RPC de identidade, testes de papel
     & Mai./Jun. & Mikael \\
  \rowcolor{rowGray}
  14 & Clerk no app: OTP telefone, sessão e SecureStore
     & Mai./Jun. & Davi \\
  15 & CI GitHub Actions, logs estruturados e \textit{tracing} (OTel)
     & Jun. 2026 & Mikael \\
  \rowcolor{rowGray}
  16 & Demandas: publicação, edição e visibilidade (contratante primeiro)
     & Jun. 2026 & Murilo, João P. \\
  17 & Descoberta: busca e filtros por serviço, município e raio
     & Jun. 2026 & Murilo, Davi \\
  \rowcolor{rowGray}
  18 & Confiança: verificação básica e portfólio com mídia
     & Jun./Jul. & Emily, João P. \\
  19 & Conexão: handoff de contato, status da demanda e avaliações
     & Jul. 2026 & Murilo, Emily \\
  \rowcolor{rowGray}
  20 & PostGIS: consultas por raio; mapas no cliente (Expo)
     & Jul. 2026 & Mikael, João P. \\
  21 & Mercado Pago: assinatura do prestador e webhooks na API
     & Jul. 2026 & Murilo, Davi \\
  \rowcolor{rowGray}
  22 & Endurecimento: testes E2E prioritários e checklist de produção
     & Jul. 2026 & Todos \\
  23 & Deploy da API (PaaS) + segredos e variáveis de ambiente
     & Jul./Ago. & Mikael \\
  \rowcolor{rowGray}
  24 & EAS Build: binários para testes internos e preparação de loja
     & Jul./Ago. & Davi \\
  25 & Testes funcionais, de integração e de segurança (app + API)
     & Ago. 2026 & Todos \\
  \rowcolor{rowGray}
  26 & Testes de usabilidade com usuários reais da AMAUC
     & Ago. 2026 & Todos \\
  27 & Correção de bugs e ajustes pós-testes
     & Ago. 2026 & Todos \\
  \rowcolor{rowGray}
  28 & Documentação final (README, relatório, diagramas) e apresentação
     & Ago. 2026 & Todos \\
\end{longtable}

% ============================================================
%  9. CONSIDERAÇÕES FINAIS
% ============================================================
\newpage
\section{Considerações Finais}

A plataforma \textbf{Opus Freelas} parte de uma necessidade concreta e
ainda pouco atendida pelo mercado: a informalidade na descoberta e na
confiança na contratação de serviços manuais e rurais na região da AMAUC.
Mais do que um exercício acadêmico, o projeto propõe resolver um problema
real com ferramentas de produção e escopo de produto explícito (fechamento
fora da plataforma na V1; assinatura do prestador como monetização inicial).

A escolha da stack --- \textbf{Expo/React Native}, \textbf{Hono}, PostgreSQL
com \textbf{PostGIS} via \textbf{Supabase}, \textbf{Clerk} e
\textbf{Mercado Pago} (assinatura) --- reflete o planejamento documentado
(GSD): uma API própria para regras e validação; banco relacional com RLS;
identidade gerenciada com JWT; cliente multiplataforma único; e pagamentos
focados no modelo de negócio da plataforma, sem intermedir o valor do
serviço entre as partes na primeira versão. Cada decisão busca maximizar
entrega e reduzir acoplamento desnecessário.

A metodologia GSD com vibe coding garante que o foco permaneça em
funcionalidade entregue, não em processo. As ferramentas de IA ---
Claude, ChatGPT, Gemini, Grok e NotebookLM --- atuam como copiloto,
acelerando o desenvolvimento sem substituir o julgamento e a revisão
crítica da equipe.

Do ponto de vista educacional, o projeto materializa a integração de todas
as competências do Curso Técnico em Informática para Internet em um produto
funcional com impacto social mensurável: trabalhadores rurais encontrando
mais clientes, contratantes encontrando profissionais confiáveis, e uma
cadeia de trabalho informal ganhando visibilidade e segurança digital.

% ============================================================
%  REFERÊNCIAS
% ============================================================
\newpage
\bibliographystyle{abntex2-alf}
\bibliography{refs}

\end{document}